\documentclass[11pt]{article}

\usepackage[margin=1in]{geometry}
\usepackage{amsmath, amssymb, amsthm}
\usepackage{mathtools}
\usepackage{hyperref}

\newtheorem{theorem}{Theorem}
\newtheorem{lemma}[theorem]{Lemma}
\newtheorem{proposition}[theorem]{Proposition}
\theoremstyle{definition}
\newtheorem{definition}[theorem]{Definition}
\newtheorem{example}[theorem]{Example}
\theoremstyle{remark}
\newtheorem{remark}[theorem]{Remark}

\title{Self-Information (Surprisal)}
\author{math-foundations / 35-self-information}
\date{}

\begin{document}
\maketitle

\subsection*{First, an example}

Before formalising surprisal, fix intuition with coins. A \emph{fair} coin has
$p(H) = p(T) = 0.5$, so its self-information in bits is
$I(H) = -\log_2(0.5) = 1$ bit --- one bit of surprise per flip, no more, no less.
\emph{Read aloud:} ``I of H equals minus log base two of one-half equals one bit.''

Now bias the coin: $p(H) = 0.9$, $p(T) = 0.1$. Then
$I(H) = -\log_2(0.9) \approx 0.152$ bits (a heads is barely surprising, since we
expected it), while $I(T) = -\log_2(0.1) \approx 3.32$ bits (a tails is genuinely
informative). Rare outcomes carry more surprisal: this is the whole picture.

Finally, flip the fair coin twice and observe $HH$. Independence gives
$I(HH) = I(H) + I(H) = 2$ bits --- surprises \emph{add} for independent events.
The rest of this lesson explains why $-\log p$ is the unique function with these
properties.

\section{Motivation}

Information theory begins with a deceptively simple question: how much
\emph{information} does the occurrence of a single event convey? Intuition
gives three desiderata.
\begin{itemize}
  \item A certain event ($p = 1$) conveys no information.
  \item Rarer events convey more information than common ones.
  \item Two independent events convey, jointly, the sum of their individual
        informations.
\end{itemize}
Shannon (1948) showed that these three properties, together with continuity,
single out a one-parameter family of functions: $I(p) = -c \log p$. The
constant $c > 0$ merely chooses units; up to that choice, $-\log p$ is the
\emph{unique} measure of surprisal.

\section{Definitions}

\begin{definition}[Self-information / surprisal]
Let $X$ be a discrete random variable with probability mass function $p$
on a countable alphabet $\mathcal{X}$, and let $x \in \mathcal{X}$ satisfy
$p(x) > 0$. The \emph{self-information} of the outcome $x$ is
\[
  I(x) \;=\; -\log p(x).
\]
\emph{Read aloud:} ``I of x equals minus log of p of x.''
The base of the logarithm fixes the units of $I$:
\[
  \log_2 \;\Rightarrow\; \text{bits}, \qquad
  \ln    \;\Rightarrow\; \text{nats}, \qquad
  \log_{10} \;\Rightarrow\; \text{hartleys.}
\]
\emph{Read aloud:} ``log base two gives bits, natural log gives nats, log base ten gives hartleys.''
\end{definition}

\begin{definition}[Self-information of an event]
For an event $A$ with $\Pr(A) > 0$, define
$I(A) = -\log \Pr(A)$. When $A = \{X = x\}$ this recovers the previous
definition.
\end{definition}

\begin{remark}
Self-information is a function of $p$, not of $x$ in any geometric sense;
two outcomes with the same probability carry the same surprisal regardless
of their identity.
\end{remark}

\section{Additivity for independent events}

\begin{theorem}[Additivity]\label{thm:additivity}
Let $A$ and $B$ be independent events with probabilities $p_A, p_B \in (0,1]$.
Then
\[
  I(A \cap B) \;=\; I(A) + I(B).
\]
\emph{Read aloud:} ``I of A intersect B equals I of A plus I of B.''
\end{theorem}

\begin{proof}
By the definition of independence,
\[
  \Pr(A \cap B) = p_A \, p_B.
\]
\emph{Read aloud:} ``probability of A intersect B equals p sub A times p sub B.''
Applying the definition of self-information,
\[
  I(A \cap B) = -\log \Pr(A \cap B) = -\log(p_A p_B).
\]
\emph{Read aloud:} ``I of A intersect B equals minus log probability of A intersect B equals minus log of p sub A times p sub B.''
Using the homomorphism property of the logarithm, $\log(xy) = \log x + \log y$
for $x, y > 0$, we obtain
\[
  -\log(p_A p_B) = -\log p_A - \log p_B = I(A) + I(B). \qedhere
\]
\emph{Read aloud:} ``minus log of p sub A times p sub B equals minus log p sub A minus log p sub B equals I of A plus I of B.''
\end{proof}

\begin{remark}
Independence is essential. If $B = A$, then $A \cap B = A$ and
$I(A \cap B) = I(A) \neq 2 I(A)$ in general. The correct
generalisation to dependent events introduces \emph{conditional}
self-information $I(B \mid A) = -\log \Pr(B \mid A)$, yielding the chain
rule $I(A \cap B) = I(A) + I(B \mid A)$.
\end{remark}

\section{Uniqueness of \texorpdfstring{$-\log p$}{-log p}}

The additivity property is, in fact, almost sufficient to characterise
surprisal.

\begin{theorem}[Shannon uniqueness, abridged]\label{thm:uniqueness}
Let $f : (0, 1] \to \mathbb{R}_{\ge 0}$ be continuous, strictly decreasing,
and satisfy
\[
  f(pq) = f(p) + f(q) \qquad \text{for all } p, q \in (0,1].
\]
\emph{Read aloud:} ``f of p times q equals f of p plus f of q for all p, q in the half-open interval zero to one.''
Then there exists a constant $c > 0$ such that $f(p) = -c \log p$.
\end{theorem}

\begin{proof}[Sketch]
Set $g(t) = f(e^{-t})$ for $t \in [0, \infty)$. The functional equation
becomes
\[
  g(s + t) = f(e^{-(s+t)}) = f(e^{-s} e^{-t})
           = f(e^{-s}) + f(e^{-t}) = g(s) + g(t),
\]
\emph{Read aloud:} ``g of s plus t equals f of e to the minus s plus t, equals f of e to the minus s times e to the minus t, equals f of e to the minus s plus f of e to the minus t, equals g of s plus g of t.''
i.e. $g$ is additive on $[0, \infty)$. Together with continuity, Cauchy's
functional equation gives $g(t) = c t$ for some constant $c$. Strict
monotonicity of $f$ forces $c > 0$. Translating back,
$f(p) = g(-\log p) = -c \log p$.
\end{proof}

\section{Worked example}

\begin{example}[Biased coin]
Let $X \in \{H, T\}$ with $\Pr(X = H) = 0.9$, $\Pr(X = T) = 0.1$. Working
in bits:
\[
  I(H) = -\log_2(0.9) \approx 0.152, \qquad
  I(T) = -\log_2(0.1) \approx 3.322.
\]
\emph{Read aloud:} ``I of H equals minus log base two of nine-tenths, approximately zero point one five two; I of T equals minus log base two of one-tenth, approximately three point three two two.''
The expected self-information is
\[
  \mathbb{E}[I(X)] = 0.9 \cdot 0.152 + 0.1 \cdot 3.322 \approx 0.469 \text{ bits},
\]
\emph{Read aloud:} ``expectation of I of X equals zero point nine times zero point one five two plus zero point one times three point three two two, approximately zero point four six nine bits.''
the Shannon entropy of $X$, which is strictly less than the maximum of $1$ bit
achieved by a fair coin.
\end{example}

\section{Connection to machine learning}

For a language model $q(x_t \mid x_{<t})$, the per-token cross-entropy
loss
\[
  \mathcal{L}_t = -\log q(x_t \mid x_{<t})
\]
\emph{Read aloud:} ``script L sub t equals minus log q of x sub t given x less than t.''
is precisely the self-information that the model assigns to the observed
token. Minimising average loss is therefore minimising mean surprisal under
the model. Perplexity, $\exp(\overline{\mathcal{L}})$, is the
exponentiated mean surprisal --- the geometric-mean inverse-probability
per token.

\end{document}
